\section{Desarrollo de \textit{Software}}

\subsection{Introducción y objetivo del código}

Una vez desarrollados los módulos y el gabinete central, se desarolla el programa de control para el microcontrolador. El software tiene como fin ser una interfaz interactiva e intuitiva entre el usuario y los distintos módulos. Cada módulo presenta su propia programación de control e interfaz. 

%% VER ESTO DE ABAJO, no me cuadra
\begin{itemize}[noitemsep]
    \item \textbf{Autodetección de módulos}, mediante una red de identificación, el código es capaz de determinar de que módulo se trata por ejemplo, convertidor Buck, recortador de onda o inversor.
    \item \textbf{Conexión segura}, utiliza un sistema de \textit{hot plugs} para verificar y asegurar la correcta conexión física del módulo, con el fin de evitar cualquier mala comunicación o falla en el mismo.
    \item \textbf{Interfaz de usuario dinámica}, Una vez que el código verifica de que módulo se trata y asegura conexión válida, el \textit{software} carga una nueva pantalla en el \textit{display}, correspondiente al módulo utilizado. Esta interfaz se diseñó para el fácil uso y aprendizaje del módulo con experimentación en tiempo real. Tiene un menú donde se puede acceder a distintas características de cada módulo y modificarlas en tiempo real, alguna de ellas pueden ser, frecuencia de conmutación, valor de \textit{setpoint}, ángulo de disparo, selección entre cargas, entre otros.
\end{itemize}

Cuando se desconecta el módulo que estaba siendo utilizado, el sistema queda en espera de una nueva conexión.

\subsection{Estructura del \textit{software}}

El código se separa en componentes, cada pieza encapsulada y con responsabilidad única. De esta forma, se simplifica la adición de nuevos módulos o funciones en el programa sin necesidad de restructurar todo el código.
% Estructura del código y configuración
% MATEO PARA ALEJO: Module connection teóricamente no es un componente aparte, sino que es parte del module manager. No sé de qué forma se podría poner para que quede bien.
\begin{figure}[H]
    \centering
        \resizebox{\textwidth}{!}{
            \begin{tikzpicture}
                % Paths, nodes and wires:
                \node[shape=rectangle, draw, line width=1pt, align=left, text width=6.077cm, inner sep=6pt] at (5.5, 2) {main.c\\ \begin{itemize}[noitemsep]
                        \item Inicializa interfaz de usuario 
                        \item Inicializa \textit{module manager}
                    \end{itemize}};
                \draw[line width=1pt] (5.5, 1.15) -- (5.5, 0);
                \node[shape=rectangle, draw, line width=1pt, align=left, text width=5cm, inner sep=6pt, minimum width=5cm, minimum height=2.7cm] at (9.125, -2.125) {Module Manager (common)\\ 
                    \begin{itemize}[noitemsep]
                        \item Identifica módulo conectado
                        \item Inicia y detiene módulo conectado
                    \end{itemize}};
                \draw[line width=1pt, -latex] (5.5, -0) -| (9.25, -0.8);
                \node[shape=rectangle, draw, line width=1pt, align=left, text width=5cm, inner sep=6pt, minimum width=5cm, minimum height=2.7cm] at (1.625, -2.125) {UI Framework (common)\\
                    \begin{itemize}[noitemsep]
                        \item Configuración de LVGL 
                        \item Configuración \textit{display}
                        \item Configuración Encoder
                    \end{itemize}};
                \draw[line width=1pt, -latex] (5.5, -0) -| (1.625, -0.8);
                \draw[line width=1pt] (9.375, -3.5) -- (9.375, -4.25);
                \draw[line width=1pt, -latex] (12.25, -4.25) -| (8.25, -4.875);
                \node[shape=rectangle, draw, line width=1pt, align=left, text width=5cm, inner sep=6pt, minimum width=5cm, minimum height=3.3cm] at (2.25, -6.5) {Módulo Dimmer\\
                    \begin{itemize}[noitemsep]
                        \item Inicia interfaz de usuario propia
                        \item Configura perisféricos para el funcionamiento
                        \item Lógica de control
                    \end{itemize}};
                \node[shape=rectangle, draw, line width=1pt, align=left, text width=5cm, inner sep=6pt, minimum width=5cm, minimum height=3.3cm] at (8.25, -6.5) {Módulo Buck\\
                    \begin{itemize}[noitemsep]
                        \item Inicia interfaz de usuario propia
                        \item Configura perisféricos para el funcionamiento
                        \item Lógica de control
                    \end{itemize}};
                \draw[line width=1pt, -latex] (14.25, -4.25) -- (14.25, -4.875);
                \node[shape=rectangle, draw, line width=1pt, align=left, text width=5cm, inner sep=6pt, minimum width=5cm, minimum height=3.3cm] at (14.25, -6.5) {Módulo Inversor\\
                    \begin{itemize}[noitemsep]
                        \item Inicia interfaz de usuario propia
                    \end{itemize}};
                \node[shape=rectangle, draw, line width=1pt, align=left, text width=5cm, inner sep=6pt, minimum width=5cm, minimum height=2.7cm] at (15.875, -2.125) {Module Connection (common)\\
                    \begin{itemize}[noitemsep]
                        \item Detecta conexión y desconexión de módulos
                        \item Espera hasta lograr conexión estable
                    \end{itemize}};
                \draw[line width=1pt, -latex] (11.83, -2.125) -- (13.15, -2.125);
                \draw[line width=1pt, -latex] (8.25, -4.25) -| (2.25, -4.875);
                \node[shape=rectangle, draw, line width=1pt, align=center, text width=5cm, inner sep=6pt, minimum width=5cm, minimum height=3.3cm] at (20.25, -6.5) {Nuevo módulo};
                \draw[line width=1pt] (12.25, -4.25) -- (16.8, -4.25);
                \draw[line width=1pt, dash pattern={on 1pt off 8pt}] (16.8, -4.25) -- (18, -4.25);
                \draw[line width=1pt, -latex] (18, -4.25) -| (20.25, -4.875);
            \end{tikzpicture}
        }
    \caption{Estructura del código implementado}
    \label{fig:soft:diagrama-arquitectura}
\end{figure}

En la \cref{fig:soft:diagrama-arquitectura} se observa un diagrama de la estructura del código. A continuación, se explica el funcionamiento de los componentes más importantes:

\subsubsection{\textit{UI Framework} - Estructura de interfaz de usuario}
Este componente se encarga de inicializar y configurar la librería LVGL (Light and Versatile Graphics Library, Librería de Gráficos Liviana y Verśatil en español), \parencite{LVGL}. Mediante esta librería, se controla el \textit{display} TFT LCD de \SI{2.8}{pulgadas}, creando las tareas, \textit{mutex}, \textit{handlers}, \textit{handles} y punteros necesarios para proveer una API (\textit{Application Programming Interface}, interfaz de programación de aplicación en español) para los módulos. También se encarga de configurar el IO (\textit{Input Output}, Entrada y Salida en español) para el funcionamiento del encoder rotativo con el que el usuario navega la interfaz en el \textit{display}.

De esta forma, la interfaz de usuario queda desacoplada del resto del sistema, permitiendo que cada módulo solo implemente la lógica y pantallas específicas sin manipular directamente drivers o detalles de sincronización.

\subsubsection{\textit{Module Manager} - Gestor de módulos}
Este componente se compone de dos programas separados, \textit{Module Connection} (Conexión de Módulo) y \textit{Module Manager}, así como un archivo de encabezado registro global de módulos, \textit{module\_registry.h}.

\paragraph{\textit{Module Connection}}
El primero se encarga de gestionar la conexión y desconexión de cada módulo, con interrupciones configuradas para ambos pines de HotPlug HP\_Power y HP\_Signal. Estos son los pines de los extremos del conector PCIe x4, que permiten detectar que la conexión del módulo haya sido exitosa, ya que sus \textit{pads} son más cortos que los demás.

Se configuró una máquina de estados para monitorear el estado de conexión: 
\begin{itemize}[noitemsep]
    \item \textit{CONNECTION\_INIT}: inicalización de programa
    \item \textit{CONNECTION\_DISCONNECTED}: módulo desconectado, señalizar a \textit{Module Manager}.
    \item \textit{CONNECTION\_DETECTING}: módulo conectado, inicando \textit{timer} de \SI{1}{\second} para verificar conexión segura y antirrebote.
    \item \textit{CONNECTION\_CONNECTED}: módulo conectado, señalizar a \textit{Module Manager}.
    \item \textit{CONNECTION\_FAULT}: falla detectada, detener máquina de estados.
\end{itemize}
El cambio de estados se realiza mediante notificaciones liberadas desde las interrupciones, despertando a una tarea general.

\paragraph{\textit{Module Manager}}
Este módulo se encarga de definir una estructura de módulo (\textit{module\_t}), identificar el módulo conectado e iniciar y detener cada módulo. Cada módulo se define en su propio archivo de control mediante una estructura (\textit{module\_t}), la cual contiene: su nombre, punteros a sus funciones de inicio y parada, valores en milivolts de identificación.

Se configuró una máquina de estados para gestionar los módulos: 
\begin{itemize}[noitemsep]
    \item \textit{MANAGER\_INIT}: inicalización de programa
    \item \textit{MANAGER\_READY}: ningún módulo corriendo, esperando conexión de módulo señalizada por \textit{Module Connection}. Una vez obtenida, notificar a tarea de identificación para detectar qué modulo se conectó, mediante lectura de \textit{ADC} (Analog to Digital Converter, Convertidor Analógico a Digital en español). Si se detecta algún módulo registrado, corrrer función de inicio de módulo.
    \item \textit{MANAGER\_RUNNING}: módulo corriendo, esperando desconexión de módulo señalizada por \textit{Module Connection}. Una vez obtenida, proceder a detener módulo corriendo actualmente.
    \item \textit{MANAGER\_STOPPING}: deteniendo módulo previamente conectado.
    \item \textit{MANAGER\_FAULT}: falla detectada, detener máquina de estados.
\end{itemize}


% Hablar sobre cada parte del código, lo más importante
% Diagrama de flujo

\subsection{Diseño de interfaz de usuario}
Para realizar una interfaz gráfica intuitiva y fácil de utilizar, se empleó el software SquareLine Studio, \parencite{squareline}. Es un editor visual de interfaz de usuario que permite crear interfaces de forma rápida y sencilla para aplicaciones integradas y de escritorio, con una fácil escalabilidad e implementación en microcontroladores. Cabe aclarar que se hizo uso de la licencia gratuita de este software. La elección de utilizar este software fue por el fácil uso. Además, de ser versatil y ampliamente utilizado en la industria.
% Explicar acá qué consideraciones se tomaron en cuenta al diseñar las pantallas, qué programa se utilizó, por qué, referenciarlo, etc.

Cada módulo cuenta con su propia interfaz lanzada cada vez que es conectado en el gabinete. Las interfaces se mantuvieron sencillas y con la información útil y necesaria; se utilizó una tipografía que facilita la lectura, y además se utilizaron colores armónicos.

\subsection{Programación de módulos}
Cada programación de módulo se compone de dos partes importantes: el control del módulo y su interfaz. A continuación se detallan los programas utilizados para los tres módulos desarrollados en el proyecto.

\subsubsection{Módulo convertidor Buck}
% Explicar máquina de estado, utilización de timers, lógica de control.
El control de este módulo fue diseñado para lograr un control preciso de la tensión de salida y, al mismo tiempo, ofrecer una interfaz de usuario fluida e interactiva. La arquitectura del software está estratégicamente dividida en dos componentes principales que operan de forma concurrente: el núcleo de control (\textit{buck\_control.c}) y el gestor de la interfaz de usuario (\textit{buck\_interface.c}). Esta separación es fundamental para aislar las operaciones críticas en tiempo, como el bucle de control, de las operaciones no críticas, como la actualización de la pantalla, garantizando así la estabilidad del sistema.

\paragraph{Arquitectura de software}
\begin{itemize}[noitemsep]
    \item \textit{buck\_control.c}: responsable de la configuración de todos los periféricos de hardware necesarios, como el convertidor analógico-digital (ADC) para la lectura de la realimentación, el generador de modulación por ancho de pulso (LEDC PWM) para controlar el interruptor de potencia y el temporizador de propósito general (GPTimer) que dicta la frecuencia del bucle de control. Además, implementa la lógica de control principal, incluyendo la máquina de estados y el algoritmo del controlador Proporcional-Integral-Derivativo (PID).
    \item \textit{buck\_interface.c}: este componente realiza la gestión de la interfaz gráfica de usuario (GUI) del módulo, utilizando la librería LVGL. Se encarga de la creación, destrucción y navegación de las distintas pantallas y elementos interactivos específicos del convertidor Buck.
\end{itemize}

\paragraph{Máquina de estados}
El comportamiento del convertidor se rige por una máquina de estados gestionada a través de la variable \textit{control\_mode}. El acceso a esta variable está protegido por un mutex (\textit{xControlModeMutex}) para evitar condiciones de carrera cuando el modo es modificado por la interfaz de usuario y leído por la tarea de control. Los estados definidos son:

\begin{itemize}[noitemsep]
    \item \textit{CONTROL\_MODE\_IDLE}, Estado seguro por defecto. La salida PWM se fuerza a cero, desactivando el convertidor.
    \item \textit{CONTROL\_MODE\_FIXED\_PWM}, Modo de control en lazo abierto. El usuario controla directamente el ciclo de trabajo del PWM a través de un \textit{slider} en la UI.
    \item \textit{CONTROL\_MODE\_PID}, Modo de control de tensión en lazo cerrado. El sistema regula activamente la tensión de salida para que coincida con el \textit{setpoint} definido por el usuario, implementando el bucle de control de realimentación completo.
\end{itemize}

\subsubsection{Módulo recortador de onda}
% Explicar máquina de estado, utilización de timers, lógica de control.

\subsubsection{Módulo inversor de onda cuadrada}
% Explicar que no tiene programa de control y que solo tiene una pantalla informativa.